\title{Secondary Education}

\section*{Mathematics}

This page contains a set of materials and resources for teaching math and programming in high school.
\begin{itemize}
  \item \href{/teaching/math-1eso/}{Matemáticas 1º ESO}
\end{itemize}

\subsection*{Resources}

\subsubsection*{Books}

\begin{itemize}
  \item Editorial Síntesis: \href{https://www.sintesis.com/matem%C3%A1ticas%3A$20cultura%20y%20aprendizaje-72}{Colección EGB} / \href{https://www.sintesis.com/educación%20matemática%20en%20secundaria-70/}{Colección ESO}
  \begin{itemize}
    \item \emph{Números y Algoritmos} (Gairín y Sancho)
    \item \emph{El problema de la medida}
    \item \emph{Estimación en el cálculo y en la medida} \textbf{¡COMPRAR!}
    \item \emph{Azar y Probabilidad}
    \item \emph{Aprender a enseñar matemáticas en la ESO} \textbf{RECOMENDADO EN JAEM}
  \end{itemize}
  \item \href{https://www.mathshell.com/}{Shell Centre}.
  \begin{itemize}
    \item \emph{Patrones y números} (Bloque de números y preálgebra)
    \item \emph{Funciones y gráficas}
  \end{itemize}
  \item \emph{Matemáticas para aprender a pensar} (Vila y Callejo)
  \item \href{http://www.ugr.es/~jgodino/edumat-maestros}{Edumat}
  \item \href{https://www.grupzero.cat/}{Grup Zero}. Materiales, aunque algo avanzados, abarcan múltiples ramas de la educación en matemáticas desde los problemas.
  \item PROBLEMAS PARA PENSAR EN MATEMÁTICAS (1º ESO). Serie de Gregorio Morales y José Ignacio Úbeda para problemas y dinámicas Thinking Classrooms en 1º ESO.
\end{itemize}

\subsubsection*{PhD Thesis}

\begin{itemize}
  \item Números enteros (Eva Cid) \href{https://web.archive.org/web/20210119170137/http://www.atd-tad.org/wp-content/uploads/2016/12/TESIS_EVA_CID-1.pdf}{pdf} \href{https://drive.google.com/file/d/0B0VFDUr_rGnYcjlSMDZfdkxscHM/view?usp=sharing}{Anexos}
  \item Números racionales (Rafael Escolano) \href{https://zaguan.unizar.es/record/84666/files/TESIS-2019-167.pdf}{pdf}
  \item Proporcionalidad (Sergio Martínez-Juste) \href{https://uvadoc.uva.es/handle/10324/52863}{pdf}
\end{itemize}

\subsubsection*{Activities and blogs}

\begin{itemize}
  \item \href{https://www.map.mathshell.org/index.php}{Mathematics Assessment Project} se pueden encontrar materiales muy interesantes para distintas ramas.
  \item \href{https://resolve.edu.au/teaching-resources}{ReSOLVE}
  \item \href{https://nrich.maths.org/maps/secondary}{nrich}
  \item \href{https://donsteward.blogspot.com}{Don Steward}
  \item \href{http://puntmat.blogspot.com}{Puntmat}
  \item \href{https://educa.aragon.es/documents/20126/2773111/%5B02.26%5D+Matem%C3%A1ticas.pdf/db9a15af-ba0a-e11e-5c13-a7c78ce6f901?t=1661254613594}{Currículo Aragón LOMLOE}. Múltiples referencias bibliográficas y propuestas de actividades SOTA.
  \item \href{https://tierradenumeros.com/}{Pablo Beltrán-Pellicer}. Especialmente interesante su \href{https://twitter.com/pbeltranp}{twitter}, y sus artículos sobre \emph{Enseñanza de la probabilidad}.
  \item \href{https://docs.google.com/spreadsheets/d/1N7LqzVwYC8UVUWKfm4ERPJ1dqVJwRHofIAMz6mS1vvo/edit#gid=0}{Actividades} de Lola Morales con Desmos. Muy interesantes.
  \item \href{https://twitter.com/SergioMJGR/status/1233098923029471232}{Unidad Didáctica 1º/2º} basada en la tesis de Escolano
  \item Modelo de barras de Singapur \href{https://sites.google.com/view/modelodebarras/home}{Curso de Belen Palop}
  \item \href{https://www.educa2.madrid.org/web/matesatulado/numeros-naturales}{Resolución de problemas de Manuel Domínguez}
  \item \href{http://docentes.educacion.navarra.es/jjimenei/index.html}{Cálculo mental}
\end{itemize}

\subsection*{Olympiads and Tournaments}

\begin{itemize}
  \item \href{https://www.ucm.es/sociedadpuigadam/olimpiada-matematica-fase-local-comunidad-de-madrid}{Olimpiada Matemática Española}. Enfocado en alumnos de bachillerato, se puede aplicar a alumnos que presentan ventajas a partir de 3º ESO, aunque si no han entrenado lo suficiente les puede resultar frustrante.
  \item \href{https://www.concursoprimavera.es/#concurso}{Concurso de Primavera}. Para despertar el interés con problemas asequibles. Perfecto para hacerles ganar confianza en sus habilidades.
\end{itemize}

\section*{Computer Science}

\subsection*{Resources}

\begin{itemize}
  \item \href{https://code.org/}{code.org} A resourceful web for the current curriculum.
  \item \href{https://tomorrowcorporation.com/}{Tomorrow Corporation} A funny approach to coding through playing.
  \item \href{https://codebreaker.teachable.com/}{Code Breaker} Some ideas for activities.
  \item \href{https://thecodingtrain.com/}{The Coding Train} Some ideas for activities. Challenges directed. Could be an option to make fortnight challenges.
  \item \href{https://www.aceptaelreto.com/}{¡Acepta el reto!}. Some problems to code. Maybe suitable for 2º ESO.
\end{itemize}
